% !TeX root = ../papers/nc-p3.tex
\section{三次曲線 $p=0$ の全点正則性}
\label{sec:rank-fifteen-cubic-zero-disk}
\begin{theorem}\label{thm:rank-fifteen-cubic-zero-disk}
標数零の形式的な階数15曲面で、$p=0$、すなわち $x^3=27y$ 上の
全ての窓は標準Koszul左右AS正則である。各始点で経路の直和は
四中心四次式の多項式環上自由階数256となる。
この曲線の正規化は $j\ne\pm1$ の直線で、
\[
 x=3(1-j^2)/4,\qquad y=(1-j^2)^3/64,\qquad q=j(j^2+3)/4 .
\]
$x=3,3/4$ の二つの点は最小周期1、残る全点は最小周期2である。
\end{theorem}

\begin{lemma}\label{lem:rank-fifteen-cubic-zero-disk}
剰余標数3の離散付値を持つ標数零の通常な窓で、
ある形式的表示が $0<\nu(H)$、$\nu(K-3)\ge25\nu(H)$ を満たすとする。
この窓は定理の正則性と自由性を持つ。
\end{lemma}
\begin{proof}
まず $K=3$、局所変数を $H$ とし、
$D=\mathbb Z_{(3)}[\sqrt2]$ 上の冪級数で考える。
$3$ はこの二次拡大で惰性を持つ。$P=s+s^{-1}$、$Q=s-s^{-1}$、
$P'=t+t^{-1}$、$Q'=t-t^{-1}$ の解を
\[
 P=\sqrt{8+2H},\quad Q=\sqrt{4+2H},\quad
 P'=\sqrt{8-2H},\quad Q'=\sqrt{4-2H},
\]
とし、四次根を $n=\sqrt{(8+\sqrt{64-4H^2})/2}$
と取る。定数項を順に $2\sqrt2,2,2\sqrt2,2,2\sqrt2$ とし、
内側の平方根は8に取る。いずれも単根である。
$N=H/n$ とし、八係数の規格化式から四つの巡回隣接像を得る。
$H=0$ は階数9のため、二次関係をそのまま特殊化すると
長さ4の余分な経路が現れる。

\path{computations/verify_rank_fifteen_cubic_zero_disk.py} は係数を
$D[H]/H^{25}$ で出発し、単元判定は定数項の二つの有理座標を3で還元して行う。
各経路行列で単元主元を消去する。目標階数に達する前に単元が尽きたら、
残った全ての係数の定数項が $\mathbb Q(\sqrt2)$ で厳密に零と確認し、
その全行を $H$ で割る。$H\ne0$ の窓ではこの行も関係である。
既知の経路次元の下界に達するまで消去するので、自由列の語は
実際の分数体上の経路基底となり、各語の座標を求められる。
長さ1から5で、$H$ による除算の回数の全始点・全文字での最大は
順に $0,0,0,2,2$ となる。従って得た経路基底での四次中心方程式に
使える精度は少なくとも21である。

各文字で中心次元 $3,2,0,0$ の下界に達するまで単元主元だけで消去できる。
よって得た自由座標から実際の中心切断を選べる。
文字零は三自由列の後二本、文字一は両自由列を取る。
各始点でこの四式を長さ4の始端に課し、以後も末尾の二次関係を挿入する。
長さ1から13で、各文字のうち最大の精度損失は全四始点とも
\[
                 0,0,0,2,4,2,2,2,2,2,2,2,0
\]
であり、最後にも一桁残る。各段で選択した消去は必要な単元主元と厳密な
定数項の零だけを使う。切断商の次元上界は $(1+t+t^2+t^3)^4$ の係数となり、
長さ13で消える。中心性による自由加群の上界、未切断経路の一般点からの下界、
定理\ref{thm:klein-rational-period-four}の左右Koszul複体の議論により主張を得る。

一般の $K=3+H^{25}T$ では $T$ は整である。
同じ定数根を使えば四次根と平方根は全て単根で、係数は
$D[T][[H]]$ で上の係数と $H^{25}$ を法として一致する。
各除算後も計算した桁数の合同は保たれる。従って証明で使った定数項は全て
$T$ に依存せず、完備付値環で特殊化して同じ上界を得る。
同じ $(x,y)$ における他の表示には基底変更・移動・反転で移せる。
\end{proof}

\begin{proof}[定理の証明]
$y$ が代数的なら数体の3の上の付値を選ぶ。$\nu(y)<0$ なら
$x^3=27y$ より定理\ref{thm:rank-fifteen-pole-cone}、
$\nu(y)=0$ なら $\bar x=0$ より定理\ref{thm:rank-fifteen-residue-line}を適用できる。
$\nu(y)>0$ では $K=3$, $H^2=16x/3$ とすると所定の $(x,y)$ を与え、
$\nu(H)=\nu(y)/6>0$ である。上の補題がこの最後の領域を尽くす。
超越的な $y$ には $\mathbb Q(y)$ の3進Gauss付値を使えば $y$ は単元である。
通常表示から外れる二点は定理\ref{thm:rank-fifteen-formal-boundary}で認証済みである。
最後に $q^2=(1-x/3)^2(1-4x/3)$ を $j=q/(1-x/3)$ で正規化すれば上式を得る。
$j^2=-3$ の二点は同じ節点に移り、$j=0$ はもう一つの $q=0$ の点に移る。
一位置移動は $(p,q)\mapsto(-p,-q)$ なので、完全な縮約不変量から最小周期も従う。
\end{proof}

\begin{corollary}
この3進の十分条件は正則性平面の二座標からも検査できる。
$\alpha=\nu(3)$, $\beta=\nu(y)>0$, $\delta=\nu(x^3/y-27)$ とする。
$\delta>6\alpha$ かつ $3\delta\ge12\alpha+25\beta$ ならば、
同じ正則性と自由性が成立する。
\end{corollary}
\begin{proof}
$\delta=+\infty$ は定理による。それ以外で $e^2=x^3/y-27$ と取る。
$D=(U+3)^3/(U-1)^2$ から $D-27=U(U-9)^2/(U-1)^2$ である。
$K=3+(e/9)z$ と置く。$\nu(e)>3\alpha$ なので方程式
$(K/3)((K+3)/3)z=K^2-1$ の還元は $2z=-1$ で単根である。
この解で $U=K^2$ とすれば所定の $D$ を与え、$U-1$ は単元である。
$V=x(U-1)^2/(U+3)$、$H^2=V$ より $\nu(H)=\beta/6$、
$\nu(K-3)=\delta/2-2\alpha$ なので補題の不等式が成立する。
\end{proof}
